% =====================================================================
\section{El problema de dos cuerpos}
\label{sec:dos-cuerpos}
% =====================================================================

Es el sistema de partículas más sencillo y el único que, en general, se resuelve
reduciéndolo a problemas de \emph{una} partícula. Notación: $M=m_{1}+m_{2}$,
\[
  \vb{r}\defeq\vb{r}_{12}=\vb{r}_{1}-\vb{r}_{2},\qquad
  \vrel\defeq\dot{\vb{r}}=\vb{v}_{1}-\vb{v}_{2},\qquad
  \Rcm=\frac{m_{1}\vb{r}_{1}+m_{2}\vb{r}_{2}}{M},\qquad \Vcm=\dot{\vb{R}}\cm .
\]

\begin{definicion}[Masa reducida]
\label{def:mu}
\begin{equation}
  \resultado{\frac1\mu\defeq\frac1{m_{1}}+\frac1{m_{2}}
  \quad\Longleftrightarrow\quad \mu=\frac{m_{1}m_{2}}{m_{1}+m_{2}}} ,
  \qquad \mu<\min(m_{1},m_{2}).
  \label{eq:mu}
\end{equation}
Si $m_{2}\gg m_{1}$, $\mu\simeq m_{1}$; si $m_{1}=m_{2}=m$, $\mu=m/2$.
\end{definicion}

\subsection{Desacoplamiento: CM y movimiento relativo}

Las ecuaciones de movimiento, con $\vb{f}_{12}=-\vb{f}_{21}$ la fuerza interna,
\begin{equation}
  m_{1}\ddot{\vb{r}}_{1}=\vb{F}_{1}\ext+\vb{f}_{12},\qquad
  m_{2}\ddot{\vb{r}}_{2}=\vb{F}_{2}\ext+\vb{f}_{21},
  \label{eq:dos-cuerpos}
\end{equation}
están acopladas: $\vb{f}_{12}$ depende de $\vb{r}_{1}$ y de $\vb{r}_{2}$.

\begin{teorema}[Desacoplamiento]
\label{teo:desacoplamiento}
Supóngase que
\begin{enumerate}[label=\textup{(\roman*)},leftmargin=2em]
  \item la fuerza exterior total sólo depende de la posición del CM:
        $\vb{F}_{1}\ext+\vb{F}_{2}\ext=\vb{F}\ext(\Rcm)$;
  \item las fuerzas exteriores producen aceleraciones cuya diferencia sólo
        depende de $\vb{r}$: $\dfrac{\vb{F}_{1}\ext}{m_{1}}-\dfrac{\vb{F}_{2}\ext}{m_{2}}=g(r)\,\vb{r}$.
\end{enumerate}
Entonces, con $\vb{f}_{12}=f(r)\,\uv{r}$ central, el problema se separa en dos
problemas independientes de una partícula:
\begin{equation}
  \resultado{M\ddot{\vb{R}}\cm=\vb{F}\ext(\Rcm),\qquad
  \mu\,\ddot{\vb{r}}=\vb{f}_{12}(\vb{r})+\mu\,g(r)\,\vb{r}} .
  \label{eq:desacoplado}
\end{equation}
En particular, si el sistema está aislado: $\Vcm=\text{const}$ y $\mu\ddot{\vb{r}}=f(r)\,\uv{r}$.
\end{teorema}

\begin{proof}
\begin{enumerate}
  \item \emph{Suma} de \eqref{eq:dos-cuerpos}: $m_{1}\ddot{\vb{r}}_{1}+m_{2}\ddot{\vb{r}}_{2}=M\ddot{\vb{R}}\cm=\vb{F}_{1}\ext+\vb{F}_{2}\ext$, pues $\vb{f}_{12}+\vb{f}_{21}=\vb{0}$. Por (i), sólo aparece $\Rcm$.
  \item \emph{Diferencia} de aceleraciones:
  \[
    \ddot{\vb{r}}=\ddot{\vb{r}}_{1}-\ddot{\vb{r}}_{2}
    =\frac{\vb{F}_{1}\ext}{m_{1}}-\frac{\vb{F}_{2}\ext}{m_{2}}
    +\Bigl(\frac1{m_{1}}+\frac1{m_{2}}\Bigr)\vb{f}_{12}
    \overset{\text{(ii)}}{=}g(r)\,\vb{r}+\frac{\vb{f}_{12}}{\mu} .
  \]
  Multiplicando por $\mu$ se obtiene la segunda ecuación de \eqref{eq:desacoplado}.\qedhere
\end{enumerate}
\end{proof}

\begin{proposicion}[Vuelta a las partículas]
\label{prop:posiciones}
Resueltos $\Rcm(t)$ y $\vb{r}(t)$,
\begin{gather}
  \resultado{\vb{r}_{1}=\Rcm+\frac{\mu}{m_{1}}\,\vb{r},\qquad
  \vb{r}_{2}=\Rcm-\frac{\mu}{m_{2}}\,\vb{r}} ,
  \label{eq:posiciones}\\
  \vb{v}_{1}=\Vcm+\frac{\mu}{m_{1}}\,\vrel,\qquad
  \vb{v}_{2}=\Vcm-\frac{\mu}{m_{2}}\,\vrel .\notag
\end{gather}
\end{proposicion}

\begin{proof}
$\Rcm+\dfrac{m_{2}}{M}(\vb{r}_{1}-\vb{r}_{2})=\dfrac{m_{1}\vb{r}_{1}+m_{2}\vb{r}_{2}+m_{2}\vb{r}_{1}-m_{2}\vb{r}_{2}}{M}=\vb{r}_{1}$,
y $\dfrac{m_{2}}{M}=\dfrac{\mu}{m_{1}}$. Análogo para $\vb{r}_{2}$; se deriva para las velocidades.
\end{proof}

El CM divide el segmento que une las partículas en razón inversa a las masas
(\cref{fig:dos-cuerpos}): $\abs{\vb{r}_{1}-\Rcm}/\abs{\vb{r}_{2}-\Rcm}=m_{2}/m_{1}$.

% ---------------------------------------------------------------------
\begin{figure}[ht]
\centering
\begin{tikzpicture}[scale=1.25]
  \coordinate (O) at (0,0);
  \coordinate (A) at (4.6,1.6);   % m1
  \coordinate (B) at (1.6,3.4);   % m2
  \coordinate (G) at (2.6,2.8);   % CM = (A+2B)/3
  \draw[aux] (B) -- (A);
  \draw[vec,cian] (O) -- (A) node[pos=0.55,below,etiq,cian,inner sep=4pt] {$\vb{r}_{1}$};
  \draw[vec,rosa] (O) -- (B) node[pos=0.5,left,etiq,rosa,inner sep=4pt] {$\vb{r}_{2}$};
  \draw[vec,texto] (O) -- (G) node[pos=0.5,below right,etiq,inner sep=2pt] {$\Rcm$};
  \draw[vec,ambar,shorten >=2.6pt] (B) -- (A);
  \node[etiq,ambar] at ($(B)!0.80!(A)+(-0.18,-0.31)$) {$\vb{r}$};
  % llaves
  \draw[decorate,decoration={brace,amplitude=4pt,raise=7pt},tenue] (B) -- (G)
    node[midway,above right=9pt and 0pt,etiq,text=tenue,inner sep=1pt] {$(\mu/m_{2})\,r$};
  \draw[decorate,decoration={brace,amplitude=4pt,raise=7pt},tenue] (G) -- (A)
    node[midway,above right=9pt and 1pt,etiq,text=tenue,inner sep=1pt] {$(\mu/m_{1})\,r$};
  % puntos
  \fill[cian] (A) circle (0.10);
  \fill[rosa] (B) circle (0.15);
  \node[punto,inner sep=1.4pt] at (G) {};
  \node[etiq,cian,anchor=west] at ($(A)+(0.14,0)$) {$m_{1}$};
  \node[etiq,rosa,anchor=south east] at ($(B)+(-0.08,0.08)$) {$m_{2}$};
  \draw[tenue,line width=0.35pt] (G) -- ++(-70:0.75) node[anchor=north,peq,inner sep=2pt] {CM};
  \node[etiq,anchor=north east] at (O) {$O$};
\end{tikzpicture}
\caption{Posiciones en el problema de dos cuerpos. El CM está sobre el segmento
que une las partículas y lo divide en razón $m_{2}:m_{1}$ (aquí $m_{2}=2m_{1}$);
$\vb{r}=\vb{r}_{1}-\vb{r}_{2}$ apunta de $m_{2}$ a $m_{1}$.}
\label{fig:dos-cuerpos}
\end{figure}
% ---------------------------------------------------------------------

\subsection{Momento lineal, angular y energía}

En el sistema del CM ($\vb{r}_{i}'=\vb{r}_{i}-\Rcm$, $\vb{v}_{i}'=\vb{v}_{i}-\Vcm$), por \eqref{eq:posiciones},
\begin{equation}
  \vb{v}_{1}'=\frac{\mu}{m_{1}}\vrel,\quad \vb{v}_{2}'=-\frac{\mu}{m_{2}}\vrel
  \quad\Longrightarrow\quad
  \resultado{\vb{p}_{1}'=\mu\,\vrel=-\vb{p}_{2}'},\qquad \vb{P}'=\vb{0} .
  \label{eq:momentos-cm}
\end{equation}
Los momentos de ambas partículas en el CM son opuestos y su módulo es el de una
partícula de masa $\mu$ con la velocidad relativa.

\begin{teorema}[König para dos cuerpos]
\label{teo:konig2}
\begin{equation}
  \resultado{\vb{L}=\Rcm\times M\Vcm+\vb{r}\times\mu\vrel},\qquad
  \resultado{K=\tfrac12MV\cm^{2}+\tfrac12\mu v_{12}^{2}} .
  \label{eq:konig2}
\end{equation}
\end{teorema}

\begin{proof}
Por König (Tema 2) basta calcular las partes internas con \eqref{eq:posiciones}:
\begin{align*}
  \vb{L}'&=\vb{r}_{1}'\times\vb{p}_{1}'+\vb{r}_{2}'\times\vb{p}_{2}'
  =\frac{\mu}{m_{1}}\vb{r}\times\mu\vrel+\Bigl(-\frac{\mu}{m_{2}}\vb{r}\Bigr)\times(-\mu\vrel)\\
  &=\mu^{2}\Bigl(\frac1{m_{1}}+\frac1{m_{2}}\Bigr)\vb{r}\times\vrel=\mu\,\vb{r}\times\vrel ,\\
  K'&=\tfrac12m_{1}\Bigl(\frac{\mu}{m_{1}}\Bigr)^{2}v_{12}^{2}+\tfrac12m_{2}\Bigl(\frac{\mu}{m_{2}}\Bigr)^{2}v_{12}^{2}\\
  &=\tfrac12\mu^{2}\Bigl(\frac1{m_{1}}+\frac1{m_{2}}\Bigr)v_{12}^{2}=\tfrac12\mu v_{12}^{2} .\qedhere
\end{align*}
\end{proof}

\begin{clave}{Idea}
El sistema equivale a \emph{dos partículas ficticias independientes}: una de masa $M$
en $\Rcm$ y otra de masa $\mu$ en $\vb{r}$. Si está aislado, la primera es libre y la
segunda describe el movimiento bajo una fuerza central $f(r)$ con centro fijo
(\cref{fig:orbitas}); las órbitas reales son copias de la relativa, escaladas por
$\mu/m_{1}$ y $-\mu/m_{2}$.
\end{clave}

% ---------------------------------------------------------------------
\begin{figure}[ht]
\centering
\begin{tikzpicture}[scale=0.95,
  declare function={rr(\t)=1.6/(1+0.6*cos(\t));}]
  % ---------------- (a) respecto del CM ----------------
  \draw[tenue!45,dash pattern=on 1.5pt off 2pt,line width=0.35pt] (-3.0,0) -- (1.75,0);
  \draw[fazul,line width=0.9pt,smooth,samples=180,domain=0:360]
    plot ({2/3*rr(\x)*cos(\x)},{2/3*rr(\x)*sin(\x)});
  \foreach \t in {80,280} \draw[fazul,line width=0.9pt,-{Stealth[length=5.5pt,width=4pt]}]
    ({2/3*rr(\t)*cos(\t)},{2/3*rr(\t)*sin(\t)}) -- ({2/3*rr(\t+3)*cos(\t+3)},{2/3*rr(\t+3)*sin(\t+3)});
  \draw[fhielo,line width=0.9pt,smooth,samples=180,domain=0:360]
    plot ({-1/3*rr(\x)*cos(\x)},{-1/3*rr(\x)*sin(\x)});
  \foreach \a/\op in {40/0.35,115/0.6,215/0.9}{
    \draw[texto,opacity=\op,line width=0.35pt]
      ({2/3*rr(\a)*cos(\a)},{2/3*rr(\a)*sin(\a)}) -- ({-1/3*rr(\a)*cos(\a)},{-1/3*rr(\a)*sin(\a)});
    \shade[ball color=fazul,opacity=\op] ({2/3*rr(\a)*cos(\a)},{2/3*rr(\a)*sin(\a)}) circle (0.075);
    \shade[ball color=fhielo,opacity=\op] ({-1/3*rr(\a)*cos(\a)},{-1/3*rr(\a)*sin(\a)}) circle (0.105);
  }
  % CM como marca de centro de masas
  \fill[texto] (0,0) circle (0.065);
  \draw[fondo,line width=0.6pt] (-0.065,0) -- (0.065,0) (0,-0.065) -- (0,0.065);
  \draw[texto,line width=0.4pt] (0,0) circle (0.065);
  \node[etiq,fazul,anchor=east] at (-2.8,0.45) {$m_{1}$};
  \node[etiq,fhielo,anchor=west] at (1.4,0.45) {$m_{2}$};
  \node[etiq,anchor=south] at (-0.65,1.95) {(a) Respecto del CM};
  % ---------------- (b) problema equivalente ----------------
  \begin{scope}[xshift=7.6cm]
  \draw[tenue!45,dash pattern=on 1.5pt off 2pt,line width=0.35pt] (-4.3,0) -- (1.3,0);
  \draw[facero,line width=0.9pt,smooth,samples=180,domain=0:360]
    plot ({rr(\x)*cos(\x)},{rr(\x)*sin(\x)});
  \foreach \t in {60,160} \draw[facero,line width=0.9pt,-{Stealth[length=5.5pt,width=4pt]}]
    ({rr(\t)*cos(\t)},{rr(\t)*sin(\t)}) -- ({rr(\t+2)*cos(\t+2)},{rr(\t+2)*sin(\t+2)});
  \fill[fhielo] (0,0) circle (0.07);
  \draw[fhielo!60,line width=0.35pt] (0,0) circle (0.13);
  \draw[vec,texto!85,line width=0.7pt,shorten >=1.5pt] (0,0) -- (-0.689,-1.892)
    node[pos=0.5,anchor=west,etiq,inner sep=3pt] {$\vb{r}$};
  \draw[vec,fazul,line width=0.7pt] (-0.689,-1.892) -- ++({0.964*0.85},{0.265*0.85})
    node[anchor=north west,etiq,fazul,inner sep=1.5pt] {$\vb{v}_{12}$};
  \shade[ball color=facero] (-0.689,-1.892) circle (0.085);
  \node[etiq,facero,anchor=south east,inner sep=3pt] at (-0.689,-1.892) {$\mu$};
  \node[peq,anchor=south,inner sep=4pt] at (0,0.1) {centro fijo};
  \node[etiq,anchor=south] at (-1.5,1.95) {(b) Problema equivalente};
  \end{scope}
\end{tikzpicture}
\caption{Sistema aislado con atracción gravitatoria y $m_{2}=2m_{1}$.
(a) Cada partícula recorre una elipse con foco en el CM ($\oplus$); en todo instante ambas
están alineadas con él (segmentos, más intensos cuanto más reciente es el instante).
(b) El movimiento relativo es el de una partícula de masa $\mu$ alrededor de un centro
fijo: la elipse (a) de $m_{1}$ es esta escalada por $\mu/m_{1}=\tfrac23$, y la de $m_{2}$ por
$-\mu/m_{2}=-\tfrac13$. La línea discontinua es el eje mayor.}
\label{fig:orbitas}
\end{figure}
% ---------------------------------------------------------------------

\begin{ejemplo}[Sistema Tierra--Luna]
\label{ej:tierra-luna}
La fuerza exterior es la atracción del Sol, de masa $M_{\odot}$, situado en el origen:
$\vb{F}_{i}\ext=-GM_{\odot}m_{i}\vb{r}_{i}/r_{i}^{3}$ ($i=T,L$). Como
$r_{TL}\ll R$, se toma $r_{T}\simeq r_{L}\simeq R=\abs{\Rcm}$:
\begin{enumerate}[label=\textup{(\roman*)},leftmargin=2em]
  \item $\vb{F}_{T}\ext+\vb{F}_{L}\ext\simeq-\dfrac{GM_{\odot}}{R^{3}}(m_{T}\vb{r}_{T}+m_{L}\vb{r}_{L})=-\dfrac{GM_{\odot}M}{R^{3}}\,\Rcm$, que sólo depende de $\Rcm$;
  \item $\dfrac{\vb{F}_{L}\ext}{m_{L}}-\dfrac{\vb{F}_{T}\ext}{m_{T}}\simeq-\dfrac{GM_{\odot}}{R^{3}}(\vb{r}_{L}-\vb{r}_{T})=-g\,\vb{r}_{TL}$, con $g\defeq GM_{\odot}/R^{3}\simeq\text{const}$.
\end{enumerate}
Por el \cref{teo:desacoplamiento},
\[
  M\ddot{\vb{R}}\cm=-\frac{GM_{\odot}M}{R^{3}}\Rcm,\qquad
  \ddot{\vb{r}}_{TL}=-\frac{G(m_{T}+m_{L})}{r_{TL}^{3}}\vb{r}_{TL}-g\,\vb{r}_{TL}.
\]
El último término es una corrección de marea:
$g\,r_{TL}^{3}/[G(m_{T}+m_{L})]=(M_{\odot}/M)(r_{TL}/R)^{3}\approx6\times10^{-3}$. Despreciándolo,
el CM orbita el Sol y la Luna orbita la Tierra como dos problemas de Kepler independientes.
\end{ejemplo}
